% Figure: the free-surface effect. A partly filled tank inside a heeled hull;
% the liquid surface stays horizontal, so a wedge of liquid migrates to the low
% side, shifting the tank's centre of gravity and lifting the body's effective
% G upward by a virtual rise i*rho_l / (V_disp*rho). Two panels: upright (left)
% and heeled (right). Schematic; no axis environment. No em-dashes.
\begin{figure}[htbp]
\centering
\begin{tikzpicture}[scale=0.95]
% ==== left panel: upright tank, liquid level horizontal ====
\begin{scope}[shift={(0,0)}]
  \draw[engblack, very thick] (-1.6,0) rectangle (1.6,1.8);
  % liquid filling lower part, horizontal surface
  \fill[engblue!25] (-1.6,0) rectangle (1.6,0.9);
  \draw[engblue, thick] (-1.6,0.9) -- (1.6,0.9);
  % tank liquid centroid g_t
  \fill[engred] (0,0.45) circle (1.7pt);
  \node[engred, font=\scriptsize, anchor=west] at (0.08,0.45) {$g_t$};
  \node[engblack, font=\scriptsize, anchor=north] at (0,-0.1) {upright};
\end{scope}
% ==== right panel: heeled tank, surface tilts, wedge migrates ====
\begin{scope}[shift={(4.6,0)}, rotate around={14:(0,0.9)}]
  \draw[engblack, very thick] (-1.6,0) rectangle (1.6,1.8);
\end{scope}
% draw the liquid surface in the heeled tank as a horizontal line (it stays level
% in space) clipped to the tank interior; approximate with a tilted-frame fill
\begin{scope}[shift={(4.6,0)}]
  % liquid body: lower trapezoid, surface horizontal in space (flat here)
  \fill[engblue!25] (-1.45,0.05) -- (1.45,0.05) -- (1.45,1.25) -- (-1.45,0.55) -- cycle;
  \draw[engblue, thick] (-1.45,0.55) -- (1.45,1.25);
  % migrated centroid g_t' to the low side
  \fill[engred] (0.45,0.55) circle (1.7pt);
  \node[engred, font=\scriptsize, anchor=west] at (0.5,0.5) {$g_t'$};
  % arrow showing the centroid travel
  \draw[engorange, ->] (0.0,0.62) -- (0.4,0.57);
  \node[engblack, font=\scriptsize, anchor=north] at (0,-0.1) {heeled};
\end{scope}
% virtual-rise annotation between the panels
\node[engamber, font=\scriptsize, align=center] at (3.1,2.25)
  {free surface raises\\effective $G$};
\draw[engamber, ->] (3.1,2.0) -- (3.1,1.5);
\end{tikzpicture}
\caption[The free-surface effect]{The free-surface effect in a partly
filled tank. Upright (left), the liquid centroid $g_t$ sits on the tank
centreline. When the hull heels (right), the liquid surface stays level in
space, so a wedge migrates to the low side and the liquid centroid moves to
$g_t'$. The migrating wedge shifts buoyancy support exactly as a slack
weight would, which raises the body's effective centre of gravity by a
virtual amount $\rho_l\, i/(\rho\, V_{\mathrm{disp}})$, with $i$ the second
moment of the tank's free surface. The metacentric height available to
resist heel is reduced by the same amount, independent of how much liquid
the tank holds. The remedy is to subdivide the tank with longitudinal
baffles, which cuts $i$ for each cell as the cube of the cell width.}
\label{fig:vol03ch10:free-surface}
\end{figure}
